This chapter collects the constructions and conventions used in both main parts of
the thesis. Rank-one \'etale local systems are the natural objects in the statement
of geometric class field theory. For the geometric and ramification-theoretic
arguments, however, it is more convenient to work with torsors: they are represented
by finite \'etale covers, admit concrete contracted products and quotients, and make
ramification accessible through extensions of local fields. 

If \(\Lambda\) is a finite commutative ring and \(G=\Lambda^\times\), the sheaf of
frames of a rank-one \(\Lambda\)-local system is a \(G\)-torsor, and conversely a
\(G\)-torsor determines an associated rank-one local system. 
We establish this
dictionary below and record its compatibility with pullback, tensor products,
symmetric powers, and ramification. Thus, the main results will be stated in the
language of local systems, while the constructions and ramification calculations
will generally be carried out in the language of torsors.

\section{Torsors and Rank-One Local Systems}
% $X$ with etale site, but say things here work for toposes in general and flat sites in general.

In what follows, we largely adhere to the treatment of torsors and group objects found in published notes of Alex Youcis \cite{YoucisTorsors}.
Let $\sC=(\cC, J)$ be a site and let $\cE=Sh(\sC)$ be the associated topos. 
Let $\mathcal{G}$ be a group object in $\mathcal{E}$. 
We denote by $\mathcal{G}\mathcal{E}$ the category of objects in $\mathcal{E}$ endowed with a left $\mathcal{G}$-action. 

\begin{definition}
    A \textbf{$\mathcal{G}$-torsor} in $\mathcal{E}$ is an object $\mathcal{P}$ of $\mathcal{G}\mathcal{E}$ satisfying the following conditions:
    \begin{enumerate}
        \item The structural morphism $\mathcal{P} \to 1$ is an epimorphism in $\mathcal{E}$ (i.e., $\mathcal{P}$ is locally non-empty).
        \item The map $\mathcal{G} \times \mathcal{P} \to \mathcal{P} \times \mathcal{P}$ defined by $(g, p) \mapsto (g \cdot p, p)$ is an isomorphism in $\mathcal{E}$ (i.e., $\mathcal{G}$ acts simply transitively on $\mathcal{P}$).
    \end{enumerate}
\end{definition}

Since $\mathcal{E}$ is the topos of sheaves on a site $\sC$, the definition can be reformulated in terms of covers. A $\mathcal{G}$-sheaf $\mathcal{P}$ is a $\mathcal{G}$-torsor if:
\begin{enumerate}
    \item For every object $X \in \mathcal{C}$, there exists a covering $\{ U_i \to X \} \in J$ such that $\mathcal{P}(U_i) \neq \emptyset$ for all $i$.
    \item For any $X \in \mathcal{C}$ where $\mathcal{P}(X)$ is non-empty, the action of $\mathcal{G}(X)$ on $\mathcal{P}(X)$ is simply transitively.
\end{enumerate}

A fundamental property of torsors is their local triviality: 
a $\mathcal{G}$-sheaf $\mathcal{P}$ is a $\mathcal{G}$-torsor if and only if it is locally isomorphic to the trivial torsor. 
Specifically, for every $X \in \mathcal{C}$, there must exist a cover $\{ U_i \to X \}$ such that the restriction $\mathcal{P}|_{U_i}$ is 
isomorphic, as a $\mathcal{G}|_{U_i}$-sheaf, to $\mathcal{G}|_{U_i}$ acting on itself by left multiplication.

A \textbf{morphism of $\mathcal{G}$-torsors} $f: \mathcal{P}_1 \to \mathcal{P}_2$ is a morphism of sheaves that is equivariant with respect to the $\mathcal{G}$-action.

\begin{prop}\label{prop:torsor_morphism_iso}
    Every morphism of $\mathcal{G}$-torsors is an isomorphism.
    Consequently, the category of $\mathcal{G}$-torsors in $\mathcal{E}$ forms a groupoid.
\end{prop}
\begin{proof}
    Locally on a cover trivializing both torsors, such a morphism is a $\mathcal{G}$-equivariant endomorphism
    of the trivial torsor $\mathcal{G}$, i.e.\ right translation by a section of $\mathcal{G}$, hence an
    isomorphism; and being an isomorphism may be checked on a cover.
\end{proof}

\begin{definition}
    We denote the groupoid of $\mathcal{G}$-torsors in $\mathcal{E}$ by $\mathbf{Tors}(\mathcal{E}, \mathcal{G})$ .  The set of isomorphism classes of $\mathcal{G}$-torsors is denoted by $\mathrm{Tors}(\mathcal{E}, \mathcal{G})$.
\end{definition}

Torsors exhibit functoriality with respect to the group object:

\begin{definition}
Let $\varphi: \mathcal{G}_1 \to \mathcal{G}_2$ be a morphism of group 
sheaves on $\sC$, and let $\mathcal{P}$ be a $\mathcal{G}_1$-torsor. We define the \textbf{contracted product} $\mathcal{G}_2 \times^{\mathcal{G}_1} \mathcal{P}$ as the quotient sheaf $\mathcal{G}_1 \backslash (\mathcal{G}_2 \times \mathcal{P})$, where $\mathcal{G}_1$ acts on the product by:
\[ g_1 \cdot (g_2, p) = (g_2 \varphi(g_1)^{-1}, g_1 \cdot p) \]

The contracted product inherits a natural left $\mathcal{G}_2$-action given on local sections by $h \cdot [g_2, p] = [h g_2, p]$, which endows it with the structure of a $\mathcal{G}_2$-torsor. 
\end{definition}

This construction yields a functor:
\[ \varphi_*: \mathbf{Tors}(\cE, \mathcal{G}_1) \to \mathbf{Tors}(\cE, \mathcal{G}_2), \quad \mathcal{P} \mapsto \mathcal{G}_2 \times^{\mathcal{G}_1} \mathcal{P} \]
On the level of isomorphism classes, $\varphi_*$ induces a map of pointed sets $\mathrm{Tors}(\cE, \mathcal{G}_1) \to \mathrm{Tors}(\cE, \mathcal{G}_2)$, sending the class of the trivial $\mathcal{G}_1$-torsor to the class of the trivial $\mathcal{G}_2$-torsor.

When $\mathcal{G}$ is a \textbf{sheaf of abelian groups} (an \textit{abelian sheaf}), the pointed set $\mathrm{Tors}(\cE, \mathcal{G})$ 
inherits the structure of an abelian group:
\begin{definition}
    Assume $\mathcal{G}$ is abelian sheaf. Let $\mathcal{P}_1$ and $\mathcal{P}_2$ be objects of $\mathbf{Tors}(\cE, \mathcal{G})$. 
    We define the \textbf{product} $\mathcal{P}_1 \otimes^{\cG} \mathcal{P}_2$ to be the quotient sheaf 
    $(\mathcal{P}_1 \times \mathcal{P}_2)/\mathcal{G}$, where $\mathcal{G}$ acts on $T$-points by having:
    \[ g \cdot (f_1, f_2) := (g f_1, g^{-1} f_2) \]
    The group object $\mathcal{G}$ then acts on the resulting quotient via its action on the presheaf quotient, which is given on classes by:
    \[ g \cdot [(f_1, f_2)] = [(g f_1, f_2)] = [(f_1, g f_2)] \]
    where the square brackets denote the class in the quotient set.\footnote{
    Equivantly, the sum is obtained as the contracted product of the $\cG \times \cG$-torsor $\cP_1 \times \cP_2$ along the multiplication map
$m: \cG \times \cG \to \cG$. }
\end{definition}

Let $[\mathcal{P}_1]$ and $[\mathcal{P}_2]$ be the isomorphism classes of $\mathcal{P}_1$ and $\mathcal{P}_2$ in $\mathrm{Tors}(\cE, \mathcal{G})$.
We define the sum $[\mathcal{P}_1] + [\mathcal{P}_2]$ to be the class $[\cP_1 \otimes^{\cG} \cP_2]$.
This structure turns $\mathrm{Tors}(\cE, \mathcal{G})$ into an abelian group, 
where the identity is the class of the trivial torsor and the inverse is obtained by the opposite action.

We now record the dictionary that explains the role of torsors in this thesis.
Let $\Lambda$ be a commutative ring object of $\mathcal{E}$ and let
$G=\Lambda^\times$ be its group of units. We write
$\mathbf{Pic}(\mathcal{E},\Lambda)$ for the groupoid of locally free
$\Lambda$-modules of rank one. For $\mathcal{L}$ in this groupoid, its
\textbf{frame torsor} is
\[
\operatorname{Fr}(\mathcal{L})
  :=\underline{\operatorname{Isom}}_\Lambda(\Lambda,\mathcal{L}).
\]
The group $G$ acts on frames by precomposition with multiplication on $\Lambda$.
Conversely, if $\mathcal{P}$ is a $G$-torsor, its \textbf{associated
$\Lambda$-module} is the contracted product
\[
\mathcal{P}\times^G\Lambda
  :=(\mathcal{P}\times\Lambda)/G,
\qquad
g\cdot(p,\lambda)=(gp,g^{-1}\lambda),
\]
where $G$ acts on $\Lambda$ by multiplication.

\begin{prop}[The torsor--local-system dictionary]
\label{prop:torsor_module_equivalence}
The functors
\[
\mathcal{P}\longmapsto\mathcal{P}\times^G\Lambda,
\qquad
\mathcal{L}\longmapsto\operatorname{Fr}(\mathcal{L})
\]
are mutually quasi-inverse symmetric monoidal equivalences
\[
(\mathbf{Tors}(\mathcal{E},G),\otimes^G,G)
\xrightarrow{\ \sim\ }
(\mathbf{Pic}(\mathcal{E},\Lambda),\otimes_\Lambda,\Lambda).
\]
These equivalences commute with pullback.
\end{prop}

\begin{proof}
Both assertions may be checked locally. On a cover on which $\mathcal{P}$ is
the trivial torsor $G$, the associated module
$\mathcal{P}\times^G\Lambda$ identifies with $\Lambda$. Conversely, on a
cover on which $\mathcal{L}$ is free of rank one,
$\operatorname{Fr}(\mathcal{L})$ identifies with $G$ acting on itself by
translation. These local identifications are compatible with the descent
data and give the two quasi-inverse equivalences. Applying the same argument
to products of local trivializations identifies contracted products with
tensor products. The construction is defined by fiber products and sheaf
quotients, and therefore commutes with pullback.
\end{proof}

When $\mathcal{E}=\operatorname{Sh}(X_{\text{\'et}})$ and $\Lambda$
is a finite constant ring, the objects of
$\mathbf{Pic}(\mathcal{E},\Lambda)$ are precisely the rank-one
$\Lambda$-local systems on $X$. We retain local-system language in the
statements of the main theorems and use the equivalent torsor language for
the geometric constructions and ramification calculations. We will pass
between the two languages using \Cref{prop:torsor_module_equivalence}.

Lastly, for any object $\cX \in \mathcal{E}$, there is a canonical identification between the slice category
$(\mathcal{G}\mathcal{E})/\cX$ and the category of group objects $\mathcal{G}(\mathcal{E}/\cX)$, where $\cX$ is viewed as having the 
trivial $\mathcal{G}$-action.

We denote by $\mathbf{Tors}(\cX, \cG)$ the category of $G$-torsors over $\cX$ in $\cG \cE_{/\cX}$.

%%% %%% %%% %%% %%% %%% %%% %%%
%       Torsors as Spaces
%%% %%% %%% %%% %%% %%% %%% %%%
\section{\'Etale Torsors, Covers, and Quotients}
Fix a scheme $X$, and let $G$ be a smooth affine $X$-group scheme. 
Denote by $X_{\text{Ét}}$ and $X_{\text{ét}}$ the big and small étale sites on $X$, respectively. 
The functor of points of $G$, which we denote by $G(\cdot) = \mathrm{Hom}_X(-, G)$, defines a group sheaf on 
$X_{\text{ét}}$.

\begin{definition} A \textit{principal $G$-bundle} (or \textit{principal homogeneous space}) for $G$ is a 
scheme $f : Y \to X$ equipped with a left $G$-action $\rho: G \times_X Y \to Y$ satisfying:
\begin{enumerate}
    \item The morphism $G \times_X Y \to Y \times_X Y$ sending $(g, y) \mapsto (gy, y)$ is an isomorphism of $X$-schemes.
    \item There exists an étale covering $\{U_i \to X\}$ such that $Y_{U_i} \cong G_{U_i}$ as $G$-schemes, 
    where $G$ acts on itself by left translation.
\end{enumerate}
\end{definition}

\begin{rem*}
    Since $G \to X$ is smooth and $Y$ is étale-locally isomorphic to $G$, 
    the morphism $f: Y \to X$ is smooth and affine. Note that $Y$ is generally not étale over $X$ unless $G$ is an 
    étale group scheme (e.g., a finite constant group).
    Similarly, if $G$ is finite then $Y \to X$ is finite.
\end{rem*}

Similarly to the case of $G$-torsors of a topos, we define a morphism of principal $G$-bundles to be a morphism of $X$-schemes commuting with the
$G$-action. We now have:

\begin{theorem}\label{theorem:torsors_are_realized_as_spaces}
    The morphism sending $Y \mapsto \mathrm{Hom}_X(-, Y)$ is an equivalence of categories from the category of principal $G$-bundles 
    to the category of $G$-torsors on $X_{\text{\'Et}}$
    Similarly, the morphism sending $Y$ to $\mathrm{Hom}_X(-, Y)$ to the category of $G$-torsors on $X_{\text{ét}}$ is an equivalence.
\end{theorem}

%% [REVIEW: unlabeled & never cited -- the big/small-site comparison is used implicitly whenever torsors are realized as schemes; consider labeling it and citing it (e.g. line ~381), or merging into the theorem above]
\begin{cor}
    There is a natural equivalence $\mathbf{Tors}(X_{\mathrm{\text{\'et}}}, G) \cong \mathbf{Tors}(X_{\text{\'Et}}, G)$ 
    inducing a bijection of pointed sets $\mathrm{Tors}(X_{\text{\'et}}, G) \xrightarrow{\cong} \mathrm{Tors}(X_{\text{\'Et}}, G)$ 
    which is an isomorphism of abelian groups if $G$ is abelian.
    And every $G$-torsor, in each of the above sites can be realized as a scheme, which is a prinicipal $G$-bundle.
\end{cor}

%%% %%% %%% %%% %%% %%% %%% %%%
%       Torsors and Cohomology
%%% %%% %%% %%% %%% %%% %%% %%%

% \subsection*{Torsors and Cohomology}
% We qoute without proof:
% \begin{theorem}
% There is a natural bijection of pointed sets     
% $\text{Tors}(\mathcal{G|_T}) \to \check{H}^1(T, \mathcal{G})$. Moreover, if $\mathcal{G}$ is an abelian group sheaf, it's an isomorphism of abelian groups.
% \end{theorem}

% The idea is that for every $[\mathcal{P}] \in \text{Tors}(T, \mathcal{G})$ 
% We 
% \begin{enumerate}
%     \item Choose a covering $\{U_i \to T\}$ such that $\mathcal{P}(U_i) \neq \emptyset$ for all $i$
%     \item choose sections $\alpha_i \in \mathcal{F}(U_i)$
% \end{enumerate}

% Then we note that for all $(i,j)$ the elements $\alpha_i |_{U_i \times_X U_j}$ and $\alpha_j |_{U_i \times_X U_j}$ differ by a unique element of $\mathcal{G}(U_{ij})$
% there exists a unique $s_{ij} \in \mathcal{G}(U_{ij})$ such that $\alpha_i |_{U_i \times_X U_j} = s_{ij}(\alpha_j |_{U_i \times_X U_j})$. 
% One can then easily see that $(s_{ij})$ defines an element of  $\check{H}^1(T, \mathcal{G})$
% which is independent of the choice of repreantiative $\cP$, choice of covering $\{U_i \to T\}$ and choice of sections.



%%% %%% %%% %%% %%% %%% %%% %%%
%       Constant Finite Group Torsors
%%% %%% %%% %%% %%% %%% %%% %%%

\subsection*{Constant Finite Group Torsors}

Let $G$ be a finite group, within this section, we denote the associated constant group scheme over $X$ by $\underline{G}$ to maintain a clear distinction between the group as a set and the group as a scheme. In subsequent sections, we shall follow the standard practice of identifying $G$ with its constant group scheme and omit the underline for brevity.

$\underline{G}$ is given by $\coprod_{g \in G} X$ with the action shuffeling the $X$'s according to multiplication.

By following the definitions, one sees that if $X$ be a connected scheme and $f : Y \to X$ is a finite Galois cover with Galois group $G$. 
Then, $f : Y \to X$ is a principal $\underline{G}$-bundle.
(Recall that a \textit{finite Galois cover} is a finite étale surjection $Y \to X$ with $Y$ connected and such that 
$G = \text{Aut}(Y/X)$ acts transitively on the geometric points of $Y$ lying over any geometric point of $X$).

On the otherhand if $f: Y \to X$ is a principal $\underline{G}$-bundle with $Y$ connected, then $Y$ is a finite Galois cover with
automorphism group $G$.

However, not all $\underline{G}$-torsors are connected. If $H\subset G$ is a proper subgroup then any connected finite \'etale cover $f: Y \to X$ with 
Galois group $H$ gives rise to a non-connected $\underline{G}$-torsor by looking at the induced $\underline{G}$-torsor $\varphi_*(Y)$ under the inclusion $\varphi: \underline{H} \to \underline{G}$.
On the otherhand, if we fix a geometric point $\overline{x} \to X$, then to give an homomorphism
$\rho \in \mathrm{Hom}_{\mathrm{cont}}(\pi_1^{\mathrm{et}}(X, x), G)$ is equivilant to give a connected pointed Galois cover
$(Y, \overline{y}) \to (X, \overline{x})$ with Galois group $H=\rho(\pi_1^{\mathrm{et}}(X, x)) \subset G$. 
Thus, pushing forward to $\underline{G}$ we get a prinicipal $\underline{G}$-bundle. The choice of a different geometric point $\overline{x'} \to X$
differ the homomorphism by an inner automorphism, thus we have:

\begin{theorem}\label{theorem:equivalence_fundamental_covers}
    Let $X$ be a connected scheme and $\overline{x}$ a geometric point of $X$. Suppose in addition that $G$ is a finite abstract group. 
    Define a map
    \begin{equation}
        \mathrm{Hom}_{\mathrm{cont}}(\pi_1^{\mathrm{et}}(X, \overline{x}), G) / \mathrm{Inn}(G) \to \mathrm{Tors}(X_{\text{\'Et}}, \underline{G}) 
    \end{equation}
by sending a homomorphism $\rho : \pi_1^{\mathrm{et}}(X, \overline{x}) \to G$ 
to the principal $\underline{G}$-bundle $\varphi_*(Y)$ where $Y$ is the principal $\underline{\rho(\pi_1^{\mathrm{et}}(X, \overline{x}))}$-bundle 
obtained above and $\varphi$ is the inclusion $\rho(\pi_1^{\mathrm{et}}(X, \overline{x})) \hookrightarrow G$. Then, the map 
 is a bijection of pointed sets where the trivial homomorphisms
  (which is the only element of its $\mathrm{Inn}(G)$-orbit) is the distinguished element of the left hand side.
\end{theorem}

If $G$ is abelian then $\mathrm{Inn}(G)$ is trivial, and we obtain:
\begin{cor}\label{cor:abelian_equivilance_fundamental_torsors}
    Let $G$ be a finite abelian group, $X$ a connected scheme, and $\overline{x}$ a geometric point of $X$. Then, the map from \Cref{theorem:equivalence_fundamental_covers}
     induces an isomorphism of abelian groups
    \begin{equation}
        \mathrm{Hom}_{\mathrm{cont.}}(\pi_1^{\text{ét}}(X, \overline{x}), G) \xrightarrow{\cong} \Tors(X_{\text{\'Et}}, \underline{G})
    \end{equation}
\end{cor}

For $G=\Lambda^\times$, the character in
\Cref{cor:abelian_equivilance_fundamental_torsors} is also the monodromy
character of the associated rank-one $\Lambda$-local system. In particular,
the torsor and the local system have the same restrictions to inertia and to
the upper ramification groups. Thus, one has ramification bounded by a given
modulus, or is tamely ramified along a given divisor, if and only if the other
does.

\subsection*{Quotients and Decomposition of Torsors}

Let $\mathcal{H}$ be a sheaf of groups acting on a sheaf $\mathcal{P}$ on a
site $\sC$. The \textbf{quotient sheaf} $\mathcal{P}/\mathcal{H}$ is the
coequalizer in $\operatorname{Sh}(\sC)$
\[
\mathcal{H}\times\mathcal{P}
\rightrightarrows
\mathcal{P}\longrightarrow\mathcal{P}/\mathcal{H},
\]
where the parallel arrows are the action and the projection.
Equivalently, it is the sheaf associated with the presheaf of orbits
$U\mapsto\mathcal{H}(U)\backslash\mathcal{P}(U)$. We use the notation
$\mathcal{P}/\mathcal{H}$ even though our torsor actions are written on the
left. In particular, the quotient always exists as a sheaf. It is
characterized by the following universal property:
for every sheaf $\mathcal{F}$, composition with
$\mathcal{P}\to\mathcal{P}/\mathcal{H}$ identifies morphisms
$\mathcal{P}/\mathcal{H}\to\mathcal{F}$ with the
$\mathcal{H}$-invariant morphisms $\mathcal{P}\to\mathcal{F}$. A quotient
sheaf need not in general be representable by a scheme; representability in
the finite-torsor situation used here is part of the next proposition.

Let $G$ be a finite abelian group, let $H\subseteq G$ be a subgroup, and let
$\pi:G\to G/H$ be the quotient homomorphism.

\begin{prop}\label{prop:quotient_torsors_finite}\label{prop:quotient_torsors}
Let $\mathcal{P}\to X$ be a $G$-torsor in $X_{\text{ét}}$. The quotient sheaf
$\mathcal{P}/H$ is represented by a finite étale $X$-scheme, and there is a
canonical isomorphism of $G/H$-torsors
\[
\mathcal{P}/H \xrightarrow{\sim}
\pi_*\mathcal{P}=(G/H)\times^G\mathcal{P}.
\]
Consequently, $\mathcal{P}\to X$ admits a canonical factorization
\[
\mathcal{P}\xrightarrow{\phi}\mathcal{P}/H\xrightarrow{\psi}X,
\]
where $\phi$ is an $H$-torsor and $\psi$ is a $G/H$-torsor. This
factorization is unique up to a unique isomorphism commuting with the
projections.
\end{prop}

\begin{proof}
The canonical morphism
\[
\mathcal{P}\longrightarrow\pi_*\mathcal{P},
\qquad p\longmapsto[1,p],
\]
is $H$-invariant because $\pi(H)=1$. The universal property of the quotient
sheaf therefore gives a canonical morphism
\[
\mathcal{P}/H\longrightarrow\pi_*\mathcal{P}.
\]
Choose an étale cover $U\to X$ that trivializes $\mathcal{P}$. Over $U$ this
morphism identifies with the identity of $(G/H)\times U$, since
\[
(G\times U)/H\cong(G/H)\times U
\qquad\text{and}\qquad
\pi_*\mathcal{P}|_U\cong(G/H)\times U.
\]
It is therefore an isomorphism of sheaves. The contracted product
$\pi_*\mathcal{P}$ is a $G/H$-torsor and, by
\Cref{theorem:torsors_are_realized_as_spaces}, is represented by a finite
étale $X$-scheme. Hence the same is true of $\mathcal{P}/H$.

Under the same local trivialization, the morphism
$\mathcal{P}\to\mathcal{P}/H$ is the product of the quotient map with the
identity,
\[
G\times U\longrightarrow(G/H)\times U,
\]
and is thus an $H$-torsor. Finally, suppose
$\mathcal{P}\to\mathcal{Q}\to X$ is another such factorization. Since an
$H$-torsor is locally split, every $H$-invariant morphism out of
$\mathcal{P}$ descends uniquely through $\mathcal{Q}$. Thus $\mathcal{Q}$
satisfies the same universal property as $\mathcal{P}/H$, which gives the
unique isomorphism compatible with the projections.
\end{proof}

\section{Symmetric Powers of Schemes and Torsors}
This section reviews the construction of quotients for schemes and torsors under finite group actions, specifically focusing on symmetric powers. 
To ensure these quotients exist as schemes, we utilize the framework of admissible actions from \cite{sga1}. 
Our treatment here closely follows the exposition in \cite{Guignard2018}
The definitions and results presented below are adapted from their work, the difference being that we ask $\Gamma$ to act on the left, to make the later construction of symmetric powers more natural. 
This foundation provides the necessary criteria for admissibility and base change required to define the symmetric powers of a scheme $X$
 and a $G$-torsor $\mathcal{P}$ over $X$.


Let $S$ be a scheme.

\begin{definition}[(\cite{sga1}, V.1.7).]
    \noindent
    \begin{itemize}
        \item Let $T$ be an object of a category $\mathcal{C}$ endowed with a right action of a group $\Gamma$. We say that \textbf{the quotient $T/\Gamma$ exists} in $\mathcal{C}$ if the covariant functor
    \[
    \begin{aligned}
    \mathcal{C} &\to \text{Sets} \\
    U &\mapsto \text{Hom}_{\mathcal{C}}(T, U)^{\Gamma}
    \end{aligned}
    \]
    is representable by an object of $\mathcal{C}$.
    \item Let $T$ be an $S$-scheme. An action of a finite group $\Gamma$ on $T$ is \textbf{admissible} if there exists an affine $\Gamma$-invariant morphism $f : T \to T'$ such that the canonical morphism $\mathcal{O}_{T'} \to f_* \mathcal{O}_T$ induces an isomorphism from $\mathcal{O}_{T'}$ to $(f_* \mathcal{O}_T)^{\Gamma}$.
\end{itemize}
\end{definition}

%% [REVIEW: never cited -- admissibility facts are used only implicitly (e.g. 'admissible quotient' in cor:sym-power-quotient discussion); consider labeling + citing, or compressing into prose with the SGA1 citation]
\begin{prop}
    The following holds:
    \begin{enumerate}
        \item \textbf{(\cite{sga1} V.1.3)}. Let $T$ be an $S$-scheme endowed with an admissible right action of a finite group $\Gamma$. If $f : T \to T'$ is an affine $\Gamma$-invariant morphism such that the canonical morphism $\mathcal{O}_{T'} \to f_* \mathcal{O}_T$ induces an isomorphism from $\mathcal{O}_{T'}$ to $(f_* \mathcal{O}_T)^{\Gamma}$, then the quotient $T/\Gamma$ exists and is isomorphic to $T'$.
        \item \textbf{(\cite{sga1}, V.1.8)}. Let $T$ be an $S$-scheme endowed with a right action of a finite group $\Gamma$. Then, the action of $\Gamma$ on $T$ is admissible if and only if $T$ is covered by $\Gamma$-invariant affine open subsets.
        \item \textbf{(\cite{sga1}, V.1.9)}. Let $T$ be an $S$-scheme endowed with an admissible right action of a finite group $\Gamma$, and let $S'$ be a flat $S$-scheme. Then, the action of $\Gamma$ on the $S'$-scheme $T \times_S S'$ is admissible, and the canonical morphism
            \[
            (T \times_S S')/\Gamma \to (T/\Gamma) \times_S S'
            \]
            is an isomorphism.
    \end{enumerate}
\end{prop}



Now let $X$ be an $S$-scheme and let $d \geq 0$ be an integer. The group ${S}_d$ of permutations of $\llbracket 1, d \rrbracket$ acts on the right on the $S$-scheme $X^{\times_S d} = X \times_S \dots \times_S X$ by the formula
\[
(x_i)_{i \in \llbracket 1, d \rrbracket} \cdot \sigma = (x_{\sigma(i)})_{i \in \llbracket 1, d \rrbracket}.
\]

\begin{prop}[\cite{Guignard2018} Proposition 2.27]
    If $X$ is a scheme, Zariski locally quasi-projective over $S$, then the right action of the symmetric group $S_d$ on the $d$-fold fiber product 
    $X^{\times_S d}$ is admissible. 
    Consequently, the quotient $\mathrm{Sym}_S^d(X) = X^{\times_S d}/S_d$ exists as a scheme over $S$.
\end{prop}

\begin{definition}
    Under the hypotheses of the proposition above, we define the \textbf{relative symmetric product} of $X$ over $S$ 
    of degree $d$ as the quotient
    $$\mathrm{Sym}_S^d(X) \coloneqq X^{\times_S d}/S_d.$$
    When the base scheme $S$ is clear from the context, we shall denote this quotient by 
    $X^{(d)}$.
\end{definition}

Guingard shows that when $X=\Spec(B)$ and $S=\Spec(A)$ then $\Sym_{S}^d(X)$ is representable by an affine $S$-scheme (See \cite{Guignard2018} Remark 2.28).

\begin{prop}[\cite{Guignard2018} Proposition 2.28]\label{prop:symmetric_power_flat_base_change}
If $X$ is flat and Zariski-locally quasi-projective over $S$, then $\text{Sym}_S^d(X)$ is flat over $S$. Moreover, for any $S$-scheme $S'$, the canonical morphism
\[
\text{Sym}_{S'}^d(X \times_S S') \to \text{Sym}_S^d(X) \times_S S'
\]
is an isomorphism.
\end{prop}

\subsection*{Geometric Properties of Symmetric Powers}

We record the geometric properties of symmetric powers that will be used in
both main parts of the thesis.

\begin{theorem}\label{theorem:int_geoint_implies_productint}
Let $X$ and $Y$ be integral schemes over a field $k$. If $X$ is geometrically
integral, then $X \times_k Y$ is integral. If both $X$ and $Y$ are
geometrically integral, then $X \times_k Y$ is geometrically integral.
\end{theorem}

\begin{prop}\label{prop:relative_symmetric_power_smooth}
Let $f:X\to S$ be a smooth, Zariski-locally quasi-projective morphism of
relative dimension one. Then the relative symmetric power
$X_S^{(d)}=\operatorname{Sym}_S^d(X)$ is smooth over $S$ of relative dimension
$d$.
\end{prop}

\begin{proof}Since $f: X \to S$ is smooth, the $d$-fold product $X_S^d$ is smooth over $S$. By \Cref{prop:symmetric_power_flat_base_change}, the quotient $X_S^{(d)}$ is flat over $S$. Smoothness is a fiberwise property for flat morphisms of finite presentation (cf. \stackstag{01V8}); thus, it suffices to verify the smoothness of the fibers.
    For every $s \in S$, the fiber of the symmetric power is given by:$$(\text{Sym}^d_S(X))_s \cong \text{Sym}^d_{\kappa(s)}(X_s)$$Consequently, we may reduce to the case where $S = \text{Spec}(k)$ for a field $k$.
    Since smoothness is preserved under base change to the algebraic closure, we may further assume $k = \bar{k}$.
    Let $z = \sum_{i=1}^{r} d_i \cdot x_i \in X^{(d)}$ be a point represented by an effective cycle of degree $d$, where the points $x_i \in X(k)$ are distinct and $\sum d_i = d$. Choose pairwise disjoint Zariski-open subsets $U_i \subset X$ such that $x_i \in U_i$. Let $W \subset X^{(d)}$ be the open subset consisting of cycles whose support is contained in $\bigcup U_i$ and which meet each $U_i$ with degree exactly $d_i$. There is a canonical isomorphism:
$$W \cong \prod_{i=1}^r U_i^{(d_i)}$$
Thus, it suffices to show that each $U_i^{(d_i)}$ is smooth of dimension $d_i$ at the point $d_i \cdot x_i$.

We may therefore restrict our attention to the "worst" case: the diagonal point $z = d \cdot x$. Let $R = \mathcal{O}_{X,x}$ be the local ring of the curve at $x$. Since $X$ is a smooth curve over $k$, $R$ is a regular local ring of dimension 1, i.e., a discrete valuation ring. Its completion is $\hat{R} \cong k[[t]]$. The completed local ring of the product $X^d$ at the point $(x, \dots, x)$ is:
$$\hat{\mathcal{O}}_{X^d, (x, \dots, x)} \cong \hat{R} \hat{\otimes}_k \dots \hat{\otimes}_k \hat{R} \cong k[[t_1, \dots, t_d]]$$
where $t_i$ is the uniformizer for the $i$-th copy of $X$. The symmetric group $S_d$ acts on this power series ring by permuting the variables
$t_i$. By the \textit{Fundamental Theorem of Symmetric Polynomials}, the ring of invariants is:
$$\left( k[[t_1, \dots, t_d]] \right)^{S^d} = k[[e_1, \dots, e_d]]$$
where $e_j$ is the $j$-th elementary symmetric power series in the variables $t_i$.

The ring $k[[e_1, \dots, e_d]]$ is a formal power series ring in $d$ variables. A fundamental theorem in commutative algebra states that a local ring is regular if and only if its completion is a formal power series ring over a field. Since $\hat{\mathcal{O}}_{X^{(d)}, z} \cong k[[e_1, \dots, e_d]]$, the local ring $\mathcal{O}_{X^{(d)}, z}$ is a regular local ring of dimension $d$. This implies that $X^{(d)}$ is smooth over $k$ of dimension $d$, completing the proof.
\end{proof}

\begin{cor}\label{cor:symmetric_power_curve_geometry}
Let $C$ be a smooth, projective, geometrically connected curve over a field
$k$. For every $d\geq 0$, the scheme $C^{(d)}$ is smooth and geometrically
integral. Moreover, every finite product of symmetric powers of $C$ is
geometrically integral.
\end{cor}

\begin{proof}
The smoothness follows from \Cref{prop:relative_symmetric_power_smooth}.
The curve $C$ is geometrically integral by \stackstag{0366}. After any field
extension, $C^d$ is integral and its image under the quotient map
$C^d\to C^{(d)}$ is irreducible. Since $C^{(d)}$ is smooth, it is geometrically
reduced, and hence geometrically integral. The assertion about finite products
now follows inductively from \Cref{theorem:int_geoint_implies_productint}.
\end{proof}


\subsection*{Symmetric Powers of Torsors}
Let $G$ be a finite abelian group, let $\cP$ be a $G$-torsor over an $S$-scheme $X$ in $S_{\text{Ét}}$. 


%% [REVIEW: never cited -- presumably the input for prop:symmetric_power_torsor (Guignard 2.32), but no proof ever invokes it; cite it there or remove]
\begin{prop}[{[\cite{sga1}]}, IX.5.8]
Assume that $\cP$ and $X$ are endowed with right actions from a finite group $\Gamma$ 
such that the morphism $\cP \to X$ is $\Gamma$-equivariant, and that the following properties hold:
\begin{enumerate}
    \item[(a)] The right $\Gamma$-action on $\cP$ commutes with the left $G$-action.
    \item[(b)] The right $\Gamma$-action on $X$ is admissible, and the quotient morphism $X \to X/\Gamma$ is finite.
    \item[(c)] For any geometric point $\bar{x}$ of $X$, the action of the stabilizer $\Gamma_{\bar{x}}$ of $\bar{x}$ in $\Gamma$ on the fiber $\cP_{\bar{x}}$ of $\cP$ at $\bar{x}$ is trivial.
\end{enumerate}
Then the action of $\Gamma$ on $\cP$ is admissible, and $\cP/\Gamma$ is a $G$-torsor over $X/\Gamma$ in $S_{\text{Ét}}$.
\end{prop}

By \Cref{theorem:torsors_are_realized_as_spaces}, $\cP$ is representable by a finite étale $X$-scheme.
For each $i \in  \llbracket 1, d \rrbracket$ let $p_i : X^{\times_S d} \to X$ be the projection on $i$-th factor, and let us consider the $G$-torsor
\[
p_1^{-1}\cP \otimes^G \cdots \otimes^G p_d^{-1}\cP = \KG \backslash \cP^{\times_S d}
\]
over $X^{\times_S d}$, where $\KG \subseteq G^d$ is the kernel of the multiplication morphism $G^d \to G$ (the dependence on $d$ is suppressed in the notation).
\footnote{
    Equivantly, the sum is obtained as the contracted product of the $G^d$-torsor $p_1^{-1}\cP_1 \times \dots \times p_d^{-1}\cP_d$ along the multiplication map
$m: G^d \to G$. }
The object $\KG \backslash \cP^{\times_S d}$ of $S_{\text{Ét}}$ is too representable by an $S$-scheme which is finite étale over $X^{\times_S d}$.
The group ${S}_d$ acts on the right on $\KG \backslash \cP^{\times_S d}$ by the formula
\[
(p_i)_{i \in \llbracket 1, d \rrbracket} \cdot \sigma = (p_{\sigma(i)})_{i \in \llbracket 1, d \rrbracket}.
\]
This action of ${S}_d$ commutes with the left action of $G$ on $\KG \backslash \cP^{\times_S d}$.

\medskip

\begin{prop}[\cite{Guignard2018} Proposition 2.32.]\label{prop:symmetric_power_torsor}
If $X$ is Zariski-locally quasi-projective on $S$, then the right action of ${S}_d$ on $\KG \backslash \cP^{\times_S d}$ is admissible,
so that the quotient $\cP^{[d]}$ of $\KG \backslash \cP^{\times_S d}$ by ${S}_d$ exists as a scheme over $S$.
Moreover, the canonical morphism $\cP^{[d]} \to \mathrm{Sym}_S^d(X)$ is a $G$-torsor, 
and the morphism
\[
p_1^{-1}\cP \otimes^G \cdots \otimes^G p_d^{-1}\cP \to r^{-1}\cP^{[d]}
\]
where $r : X^{\times_S d} \to \mathrm{Sym}_S^d(X)$ is the canonical projection, is an isomorphism of $G$-torsors over $X^{\times_S d}$.
\end{prop}

\begin{definition}[Symmetric powers of rank-one local systems]
\label{definition:symmetric_power_local_system}
Let $\Lambda$ be a finite commutative ring, let $G=\Lambda^\times$, and let
$\mathcal{F}$ be a rank-one $\Lambda$-local system on $X$. If
$\mathcal{P}=\operatorname{Fr}(\mathcal{F})$, we define
\[
\mathcal{F}^{[d]}:=\mathcal{P}^{[d]}\times^G\Lambda
\]
on $X^{(d)}$. Equivalently, $\mathcal{P}^{[d]}$ is the frame torsor of
$\mathcal{F}^{[d]}$. By \Cref{prop:symmetric_power_torsor} and the monoidal
compatibility in \Cref{prop:torsor_module_equivalence}, its pullback along
$r:X^{\times_S d}\to X^{(d)}$ satisfies
\[
r^{-1}\mathcal{F}^{[d]}
\cong p_1^{-1}\mathcal{F}\otimes_\Lambda\cdots\otimes_\Lambda
p_d^{-1}\mathcal{F}.
\]
Thus this definition agrees with the usual descent construction of the
symmetric power of a rank-one local system.
\end{definition}
%the $S_d$-action on $\KG \backslash \cP^{\times_S d}$ is compatible with the $S_d$-action on $X^{\times_S d}$, so that we get a canonical morphism $\cP^{[d]} \to \mathrm{Sym}_S^d(X)$.

% %Note that there is a canonical map $\cP^{\times_S d}/S_d \to \cP^$, and the $S_d$-action on $\KG \backslash \cP^{\times_S d}$ is compatible with the $S_d$-action on $X^{\times_S d}$, so that we get a canonical morphism $\cP^{[d]} \to \mathrm{Sym}_S^d(X)$.
% \textcolor{red}{This entire corollary is not exact, replace with: 
% https://gemini.google.com/app/8d257f0b8e6f5927
% or
% https://gemini.google.com/app/f36ed1280d780bea
% }

When $G$ is abelian, and letting the $S_d$-action be a left action (via the inverse), then
the condition  $\prod_{i=1}^d g_i = 1$ defining $\KG$ is invariant under permuting the coordinates.
Therefore, the left action of $S_d$ normalizes the action of $\KG$.
The group generated by these two actions is the semidirect product $\Xi = \KG \rtimes S_d$,
which is a finite group acting on the left on $\cP^d$.
By \Cref{prop:symmetric_power_torsor}, the induced action of $S_d$ on
$\KG\backslash\cP^{\times_S d}$ is admissible. Hence the iterated quotient
\[
(\KG\backslash\cP^{\times_S d})/S_d
\]
exists as an $S$-scheme. Moreover, since $\KG$ is normal in
$\Xi=\KG\rtimes S_d$, a morphism from $\cP^{\times_S d}$ is
$\Xi$-invariant if and only if it is $\KG$-invariant and the induced
morphism from $\KG\backslash\cP^{\times_S d}$ is $S_d$-invariant.
Consequently, the universal properties of the quotients give a canonical
isomorphism
\[
\Xi\backslash\cP^{\times_S d}
\cong
(\KG\backslash\cP^{\times_S d})/S_d.
\]
Here $\Xi$ acts on $\cP^{\times_S d}$ by the formula
$$((g_1, \dots, g_d), \sigma) \star (p_1, \dots, p_d) = (g_1 p_{\sigma^{-1}(1)}, g_2 p_{\sigma^{-1}(2)}, \dots, g_d p_{\sigma^{-1}(d)})$$


\begin{cor}\label{cor:sym-power-quotient}
    Under the assumptions of the above discussion, there is a natural canonical morphism 
    $$q: S_d \backslash \cP^{\times_S d} \to \Xi \backslash \cP^{\times_S d}$$
    or equivlantly
    $$ q: S_d \backslash \cP^{\times_S d} \to (\KG \backslash \cP^{\times_S d}) / S_d$$

    between the symmetric power of $\cP$ as a scheme and the symmetric power of $\cP$ as a $G$-torsor.
    We make it a convention to denote the symmetric power of $\cP$ as a scheme by 
    $\cP^{(d)}$, and the symmetric power of $\cP$ as a $G$-torsor by $\cP^{[d]}$. 
    So that the above qoutient map is also written as:
    $q: \cP^{(d)} \to \cP^{[d]}$.
\end{cor}

\begin{proof}
    $S_d$ is a subgroup of $\Xi$,
    hence any map that is $\Xi$-invariant is automatically $S_d$-invariant.
\end{proof}

\begin{rem*}
    Note that the above discussion does not require that $\cP$ is a $G$-torsor over $X$, only that it is a scheme
    with a left $G$-action.
\end{rem*}




%
% COMPATIBILITY WITH ADDITION
%
\subsection*{Additivity of Symmetric Powers}
The construction of symmetric powers is compatible with a natural addition law.
For any integers $d_1, d_2 \geq 0$, there exists a canonical \textbf{addition morphism}
\[ +_{d_1, d_2}: X^{(d_1)} \times_S X^{(d_2)} \to X^{(d_1+d_2)} \]
induced by the equivariant isomorphism $X^{\times_S d_1} \times_S X^{\times_S d_2} \cong X^{\times_S (d_1+d_2)}$ with respect to the inclusion of the product of symmetric groups $S_{d_1} \times S_{d_2} \subseteq S_{d_1+d_2}$.

\begin{prop}\label{prop:symmetric_powers_on_curves}
    Let $\cP$ be a $G$-torsor over $X$. There is a canonical isomorphism of $G$-torsors over $X^{(d_1)} \times_S X^{(d_2)}$:
    \[
    +_{d_1, d_2}^{-1} \left( \cP^{[d_1+d_2]} \right) \cong r_1^{-1} \cP^{[d_1]} \otimes^G r_2^{-1} \cP^{[d_2]}.
    \]
\end{prop}
\begin{proof}
Let $\pi = r_{d_1} \times r_{d_2}$. By the commutativity of the following diagram
\[
\begin{tikzcd}
X^{\times_S d_1} \times_S X^{\times_S d_2} \arrow[r, "\sim"] \arrow[d, "\pi"'] & X^{\times_S (d_1+d_2)} \arrow[d, "r_{d_1+d_2}"] \\
X^{(d_1)} \times_S X^{(d_2)} \arrow[r, "+_{d_1, d_2}"] & X^{(d_1+d_2)}
\end{tikzcd}
\]
and applying \Cref{prop:symmetric_power_torsor}, we obtain canonical isomorphisms:
\[
\pi^{-1} \left( +_{d_1, d_2}^{-1} \cP^{[d_1+d_2]} \right) \cong r_{d_1+d_2}^{-1} \cP^{[d_1+d_2]} \cong \sideset{}{^G}\bigotimes_{i=1}^{d_1+d_2} p_i^{-1} \cP
\]
and
\[
\pi^{-1} (r_1^{-1} \cP^{[d_1]} \otimes^G r_2^{-1} \cP^{[d_2]}) \cong \left( \sideset{}{^G}\bigotimes_{i=1}^{d_1} p_i^{-1} \cP \right) \otimes^G \left( \sideset{}{^G}\bigotimes_{j=d_1+1}^{d_1+d_2} p_j^{-1} \cP \right).
\]
Since both pullbacks identify with
\[
\sideset{}{^G}\bigotimes_{k=1}^{d_1+d_2} p_k^{-1} \cP
\]
in an $S_{d_1} \times S_{d_2}$-equivariant manner, the isomorphism descends to the quotient $X^{(d_1)} \times_S X^{(d_2)}$ by \Cref{prop:symmetric_power_flat_base_change}.
\end{proof}

\section{Basic Geometric Lemmas}

We conclude with two elementary facts about codimension-one points and
blowups. They are placed here because they enter the geometric arguments in
both later chapters.

\begin{prop}\label{prop:finite_flat_prime_divisor}
Let $f:X\to Y$ be a finite flat morphism between integral schemes of finite
type over a field $k$. If $Z\subset X$ is a prime divisor with generic point
$\eta_Z$, then $f(Z)\subset Y$ is a prime divisor, and its generic point is
$f(\eta_Z)$.
\end{prop}

\begin{proof}
    $f$ is finite hence proper hence closed so $f(Z)$ is closed subset of $Y$, it is irreducible as the image of an irreducible.
    Since $Z = \overline{\{ \eta_Z \}}$ we get:
    $$\{f(\eta_Z)\} \subset f(Z) = f(\overline{\{\eta_Z\}}) \subseteq \overline{f(\{\eta_Z\})} = \overline{\{f(\eta_Z)\}}$$
    And since $f(Z)$ is closed we get $f(Z) = \overline{\{f(\eta_Z)\}}$.

    For flat map of integral schemes we have for every $x \in X$, $y=f(x)$ the dimension formula:
    $$\text{dim}(\mathcal{O}_{X,x}) = \text{dim}(\mathcal{O}_{Y,y}) + \text{dim}(\mathcal{O}_{X_y, x})$$
    And since $\text{dim}(\mathcal{O}_{X_y, x}) = 0$ we get $\text{dim}(\mathcal{O}_{X,x}) = \text{dim}(\mathcal{O}_{Y,y})$
    concluding that $f(Z)$ is a prime divisor as well.
\end{proof}

\subsection*{Blowups}

\begin{theorem}[\stackstag{0805}]\label{theorem:blowup_flat_base_change}
Let $X_1\to X_2$ be a flat morphism of schemes, let $Z_2\subset X_2$ be a
closed subscheme, and let $Z_1$ be its inverse image in $X_1$. If $X'_i$ is
the blowup of $X_i$ along $Z_i$, then the diagram
\[
\xymatrix{
X'_1 \ar[r] \ar[d] & X'_2 \ar[d] \\
X_1 \ar[r] & X_2
}
\]
is cartesian.
\end{theorem}

\begin{prop}\label{prop:exceptional_divisor_smooth_point}
Let $X$ be a smooth $k$-scheme of pure dimension $d\geq 1$, let $z\in X(k)$, and
let
\[
\pi:\widetilde X=\operatorname{Bl}_z(X)\longrightarrow X
\]
be the blowup at $z$. Then the exceptional divisor
$E=\pi^{-1}(z)$ is isomorphic to $\mathbb P_k^{d-1}$. In particular, $E$ is
irreducible and has codimension $1$ in $\widetilde X$. Consequently, it has
a unique generic point $\eta_E$, and $\overline{\{\eta_E\}}=E$.
\end{prop}

\begin{proof}
Because $X$ is smooth of dimension $d$ at the $k$-rational point $z$, after
replacing $X$ by a sufficiently small open neighborhood of $z$ there is an
étale morphism
\[
f:X\longrightarrow \mathbb A_k^d
\]
that sends $z$ to the origin $0$. We may shrink this neighborhood further so
that $z$ is the scheme-theoretic inverse image of $0$. Since an étale
morphism is flat, \Cref{theorem:blowup_flat_base_change} reduces the assertion
to the blowup of $\mathbb A_k^d$ at the origin.

Write $x_1,\ldots,x_d$ for the affine coordinates on $\mathbb A_k^d$ and
$[u_1:\cdots:u_d]$ for the homogeneous coordinates on
$\mathbb P_k^{d-1}$. The blowup of the origin is the closed subscheme
\[
\operatorname{Bl}_0(\mathbb A_k^d)
=
\left\{
((x_1,\ldots,x_d),[u_1:\cdots:u_d])
\ \middle|\
x_i u_j=x_j u_i\text{ for all }i,j
\right\}
\subset \mathbb A_k^d\times_k\mathbb P_k^{d-1}.
\]
Over the origin all the $x_i$ vanish, so the displayed equations impose no
condition on $[u_1:\cdots:u_d]$. Hence the exceptional fiber is
\[
\{0\}\times_k\mathbb P_k^{d-1}\cong\mathbb P_k^{d-1}.
\]
It is therefore irreducible of dimension $d-1$. The blowup has dimension
$d$ (for example, each standard chart $u_i\neq0$ has $d$ affine
coordinates), so the exceptional divisor has codimension $1$.
\end{proof}

\begin{theorem}[{\stackstag{01OF}, Lemma 31.33.9}]\label{theorem:blowup_integral}
If $X$ is integral and $Z\subsetneq X$ is a proper closed subscheme, then the
blowup $\operatorname{Bl}_Z(X)$ is integral.
\end{theorem}
